\section{CNN ResNet18}
\label{chap:cnn_enfoque}

\subsection{Arquitectura y entrenamiento}
\label{sec:cnn_entrenamiento}

Las dos variantes utilizan ResNet18 preentrenada con ImageNet, una colección de imágenes de uso general \cite{he2016}. Sus capas convolucionales extraen las características visuales y la última capa las transforma en las salidas de la tarea. En el simulador, esa capa se sustituye por tres salidas, correspondientes a RBBB, elevación del ST e inversión de T. Con los datos cedidos por el HUCA se utiliza una única salida, asociada a la etiqueta de Brugada. En ambos casos se ajustan todos los pesos, tanto del extractor como de la nueva capa de clasificación. La Figura~\ref{fig:cnn_convolutional_pipeline} representa el entrenamiento y la predicción para el caso del simulador.

\begin{figure}[H]
  \centering
  % Fuente editable: assets/results/cnn_convolutional_pipeline.svg.
  \safeincludegraphics[width=0.96\textwidth]{results/cnn_convolutional_pipeline.pdf}
  \caption{Entrenamiento supervisado y evaluación de la CNN entrenada con el simulador.}
  \label{fig:cnn_convolutional_pipeline}
\end{figure}

Para presentar las imágenes a la red, se convierten a RGB, con tres canales de color, y se redimensionan a 224 por 224 píxeles. Esta conversión adapta el formato a la entrada de ResNet18 y mantiene el mismo trazado en los tres canales. A continuación se normalizan los píxeles con la media y la desviación calculadas únicamente sobre las imágenes de entrenamiento. Esos mismos valores se aplican a validación y prueba, por lo que los pacientes reservados no intervienen en la preparación de la entrada.

\begin{samepage}
El entrenamiento busca que las salidas se aproximen a las etiquetas disponibles mediante la entropía cruzada binaria. En el simulador se calcula por separado para los tres rasgos, que pueden coexistir en una imagen, mientras que en los registros clínicos se aplica a su única etiqueta. Para compensar diferencias en su frecuencia, cada salida incorpora un peso positivo igual al número de negativos dividido entre el de positivos del entrenamiento utilizado. Así se aumenta la penalización de los errores en positivos cuando estos son menos frecuentes. Con la selección balanceada de los datos cedidos por el HUCA ese peso vale uno, porque ya se utiliza el mismo número de pacientes de cada clase.
\par
\end{samepage}

Para reducir esta pérdida se emplea AdamW \cite{loshchilov2019}, que actualiza los pesos después de procesar cada lote de imágenes con una magnitud de ajuste regulada por la tasa de aprendizaje. Al recorrer todos los lotes se completa una época, tras la cual comienza un nuevo recorrido del conjunto de entrenamiento. La tasa de aprendizaje, el tamaño del lote y los demás parámetros utilizados se recogen en la Tabla~\ref{tab:config_cnn_appendix}.

\subsection{Predicción y selección del umbral}
\label{sec:umbral_cnn}

Cada salida de la red es una puntuación, o \textit{logit}, que la sigmoide transforma en una probabilidad entre cero y uno. En la CNN entrenada con el simulador se aplica por separado a las tres salidas, correspondientes a RBBB, elevación del ST e inversión de T, mientras que en la entrenada con los datos clínicos se aplica a la única puntuación asociada a Brugada, indicada con el subíndice \(B\). Ambas operaciones se expresan como
\begin{equation}
  \begin{gathered}
    p_j = \frac{1}{1+\exp(-z_j)},\quad j=1,2,3\\
    \text{Simulador}
  \end{gathered}
  \qquad
  \begin{gathered}
    p_B = \frac{1}{1+\exp(-z_B)}\\
    \text{Datos clínicos}
  \end{gathered}
  \label{eq:probabilidades_cnn}
\end{equation}
Los tres rasgos del simulador pueden coexistir, por lo que sus probabilidades no tienen que sumar uno. La clasificación positiva derivada se obtiene cuando las tres alcanzan el umbral, mientras que en los registros cedidos por el HUCA se compara con él la probabilidad \(p_B\).

El umbral determina a partir de qué probabilidad se asigna una etiqueta positiva y se elige buscando la mayor BA en los pacientes de validación. En la CNN entrenada con el simulador, la regla conjunta se evalúa a partir de la menor de sus tres probabilidades, pues las tres deben alcanzar el mismo umbral.

\begingroup
\widowpenalty=10000
En la CNN entrenada con los datos clínicos, la elección del umbral se integra en el seguimiento del entrenamiento. Al terminar cada época se predicen los ocho pacientes de validación definidos en la sección~\ref{sec:validacion_comun} y se obtiene la mejor BA con ese conjunto. Este resultado permite conservar los pesos que mejor clasifican a pacientes distintos de los utilizados para actualizarlos. El \textit{early stopping} interrumpe el ajuste cuando transcurren cinco épocas sin mejora, con un máximo de veinte, y se recuperan los pesos seleccionados para efectuar la predicción final.
\par
\endgroup

Con esos pesos se fija el umbral a partir de los mismos pacientes de validación y se clasifica al paciente reservado para prueba. De este modo, los resultados agregados se construyen con las predicciones de prueba de cada pliegue.

\subsection{Adaptación de dominio}
\label{sec:adaptacion_cnn}

Se estudian cuatro variantes para aprovechar las imágenes cedidas por el HUCA al trasladar lo aprendido con el simulador, manteniendo la predicción de los tres rasgos. La primera utiliza \textbf{SimCLR} para preparar el extractor antes del ajuste supervisado. A partir de dos versiones aumentadas de cada imagen clínica, aprende representaciones que toleren cambios de brillo, pequeños desplazamientos, ruido y desenfoque. Los pesos obtenidos inicializan después la CNN que se entrena con las etiquetas del simulador.

\textbf{CORAL} y \textbf{MMD} incorporan la adaptación durante ese entrenamiento, combinando lotes etiquetados del simulador con imágenes clínicas sin etiquetas. Mientras la pérdida supervisada aprende a distinguir los rasgos, CORAL aproxima las covarianzas de las características de ambos conjuntos y MMD reduce su discrepancia mediante kernels de base radial a varias escalas. La contribución de estos objetivos aumenta progresivamente a lo largo del ajuste.

\textbf{DANN} utiliza también ambos conjuntos, pero añade un clasificador que intenta reconocer el origen de cada imagen. La inversión de gradiente orienta al extractor a dificultar esa distinción a la vez que aprende los rasgos del simulador. Esta rama auxiliar interviene durante el entrenamiento y se retira para predecir, conservando las tres salidas de la tarea.

En las cuatro variantes, las imágenes del conjunto clínico se comparten entre pliegues, incluidas las que después se evalúan, sin utilizar sus etiquetas para entrenar. Así se estudia el aprovechamiento de su apariencia, frente al ajuste supervisado con pacientes etiquetados descrito anteriormente. La Tabla~\ref{tab:config_domain_appendix} reúne los parámetros de cada método y las condiciones de uso de los datos.

\subsection{Inspección con Grad-CAM}

La inspección con Grad-CAM complementa las métricas con una representación de las regiones relacionadas con una salida de la CNN. En este trabajo se aplica al último bloque convolucional de ResNet18, antes de reducir las características a un vector para clasificar. El mapa se calcula a partir de los gradientes de la primera puntuación de salida y se amplía al tamaño de la entrada. En la CNN entrenada con el simulador, esa salida corresponde a RBBB, por lo que el ejemplo mostrado examina ese rasgo.

\begin{samepage}
La Figura~\ref{fig:gradcam_sintetico_modelos} presenta juntos el trazado V1, el mapa normalizado y su superposición. Esta disposición permite localizar las zonas resaltadas sobre el ECG y observar si se sitúan en torno al QRS, relacionado con el rasgo representado, o en regiones periféricas como el fondo y los márgenes, indicando el color la contribución relativa dentro de esa imagen. La inspección se utiliza como apoyo cualitativo a la lectura del comportamiento del modelo y no interviene en la selección del umbral ni en el cálculo de las métricas. El fundamento del método y sus límites interpretativos se desarrollan en la sección~\ref{sec:explicabilidad}.
\par
\end{samepage}

\begin{figure}[H]
  \centering
  \safeincludegraphics[width=0.92\textwidth]{results/cnn_simulator_qrs_gradcam_example.pdf}
  \caption{Inspección de la salida RBBB de la CNN entrenada con el simulador mediante Grad-CAM. Se muestran la imagen V1, el mapa normalizado y su superposición.}
  \label{fig:gradcam_sintetico_modelos}
\end{figure}
